\documentclass[12pt]{article}
\usepackage[utf8]{inputenc}
\usepackage[T1]{fontenc}
\usepackage{hyperref}
\title{Yanom\"am\"o Language Resources for K-12 Curriculum Development}
\author{Compiled by Claude Sonnet 5.5 (claude-sonnet-5-5) for Andr\'as Kornai}
\date{}
\begin{document}
\maketitle

\section{Introduction}
Yanom\"am\"o (ISO 639-3 guu) is a Yanomaman language of Venezuela and Brazil. Web searches found no published Yanom\"am\"o primer or dictionary under that name. What exists is mostly academic work on literacy and schooling in closely related Yanomami varieties (Parima B, Sanöma, Ninam), and the relation of each variety to guu should be checked before use.

\section{Primary Educational Resources}
\begin{itemize}
\item Thesis on a syllabic literacy method adapted for the Parima B dialect of Yanomami, based on New Tribes Mission materials from the 1990s (primers, cards, teacher manuals judged appropriate): \url{https://saber.ucv.ve/jspui/handle/10872/18057}
\item Paper on a first-stage literacy book in Sanöma (a Yanomami language, so not necessarily guu), written jointly by linguists and anthropologists: \url{https://repositorio.ufmg.br/handle/1843/62336}
\end{itemize}

\section{Linguistic and Reference Materials}
\begin{itemize}
\item NSF-funded project to finish a dictionary of Ninam (a Yanomaman language), begun by Ernesto Migliazza, with scholarly and community versions: \url{https://grantome.com/grant/NSF/BCS-0963595}
\item News report mentioning a bilingual Spanish--Yanomami mini-dictionary and a flora and fauna glossary: \url{https://www.ipsnews.net/?p=13153}
\end{itemize}

\section{Religious and Community Resources}
No religious or community-organisation resources were identified in the searches.

\section{Digital and Audio Resources}
No digital or audio resources were identified in the searches.

\section{Curriculum Development Context}
A UFPB dissertation on Yanomami schooling in Maturacá (Brazil) reports Yanomami as L1 and Portuguese as L2: \url{https://repositorio.ufpb.br/jspui/handle/tede/6455}. Venezuela's 1999 constitution guarantees an intercultural bilingual education system for indigenous peoples (mentioned in the IPS report above).
\end{document}